\setcounter{section}{17}

  \zwischenueberschrift{Soal latihan}

\inputaufgabe
{}
{

Misalkan
\maabbeledisp {\varphi} {\R^n } {\R
} { \left(  x_1   , \,   \ldots  , \,   x_n                 \right) } {a_1x_1 + \cdots + a_nx_n
} {,}
suatu
\definitionsverweis {bentuk linear}{}{.}
Misalkan $M$ adalah
\definitionsverweis {grafik}{}{}
fungsi ini, yang kita pandang sebagai
\definitionsverweis {manifold Riemann}{}{.}
Tunjukkan bahwa terdapat hubungan konstan antara volume himpunan-himpunan bagian yang saling bersesuaian di $\R^n$ dan pada grafik tersebut.

}
{} {}

\inputaufgabegibtloesung
{}
{

Hitung luas permukaan bola dengan menggunakan
Korolari 17.1.

}
{} {}

\inputaufgabe
{}
{

Bahas
\definitionsverweis {permukaan putar}{}{} $S$ yang diperoleh dengan memutar

\mathdisp {M= \left\{ ( \sin  { \frac{   1  }{  y } } ,y)  \mid   y> 0        \right\}} {  }

mengelilingi sumbu $x$, yaitu $A$. Apakah $S$ merupakan
\definitionsverweis {submanifold tertutup}{}{} dari
\mathl{\R^3 \setminus A}{?} Apakah himpunan $S$
\definitionsverweis {tertutup}{}{} di $\R^3$? Apakah
\definitionsverweis {penutupan}{}{} $S$ di $\R^3$ merupakan suatu manifold?

}
{} {}

\inputaufgabegibtloesung
{}
{

Tunjukkan bahwa luas permukaan putar yang diperoleh dengan memutar grafik

\mathdisp {\Gamma = \left\{ (x,e^x)  \mid   x \leq 0        \right\}} {  }

mengelilingi sumbu $x$ lebih kecil daripada $10$.

}
{} {}

\inputaufgabe
{}
{

Pastikan bahwa
\definitionsverweis {pemetaan-pemetaan}{}{}
yang diberikan dalam
Contoh 17.6,
Contoh 17.7,
dan
Contoh 17.8
memiliki
\definitionsverweis {citra}{}{}
yang terletak pada
\definitionsverweis {sfer satuan}{}{}
dan bersifat surjektif kecuali pada suatu
\definitionsverweis {himpunan nol}{}{.}

}
{} {}

\inputaufgabe
{}
{

Tentukan
\zusatzklammer {yang terdefinisi secara parsial} {} {}
\definitionsverweis {pemetaan-pemetaan invers}{}{}
dari
\definitionsverweis {pemetaan-pemetaan}{}{}
yang diberikan dalam
Contoh 17.6,
Contoh 17.7,
dan
Contoh 17.8.

}
{} {}

\inputaufgabe
{}
{

Tunjukkan bahwa
\definitionsverweis {lingkaran bujur}{}{}
dan
\definitionsverweis {lingkaran lintang}{}{}
pada bola bumi saling
\definitionsverweis {tegak lurus}{}{.}

}
{} {}

\inputaufgabe
{}
{

Berapakah panjang
\definitionsverweis {lingkaran lintang}{}{} pada lintang $30^\circ$ di Bumi
\zusatzklammer {gunakan $6370$ km sebagai jari-jari bumi} {} {?}

}
{} {}

\inputaufgabe
{}
{

Berapa persen permukaan Bumi yang pada suatu saat disinari Matahari
\zusatzklammer {karena jaraknya sangat jauh, sinar Matahari dianggap saling sejajar} {} {?}

}
{} {}

\inputaufgabe
{}
{

Tinjau elipsoid

\mathdisp {M= \left\{ (x,y,z)\in \R^3  \mid   \ \frac{x^2}{a^2}+\frac{y^2}{b^2}+\frac{z^2}{c^2} =1         \right\}} { , }

untuk $a,b,c\in\mathbb{R}_{+}$. Hitung luas permukaan M untuk
\mathl{a= b}{.} Apa yang terjadi ketika
\mathl{a\rightarrow c}{?}

}
{} {}

\inputaufgabe
{}
{

Tentukan
\definitionsverweis {infimum}{}{}
dan
\definitionsverweis {supremum}{}{}
dari
\definitionsverweis {panjang}{}{}
citra
\definitionsverweis {lingkaran-lingkaran besar}{}{}
pada peta yang diuraikan dalam
Contoh 17.6.

}
{} {}

\inputaufgabe
{}
{

Hubungkan
Korolari 17.1
dengan
Contoh 15.10
dengan menggunakan
Latihan 16.10.

}
{} {}

\inputaufgabe
{}
{

Misalkan
\maabbdisp {\alpha} {S^2 \setminus \{N\}} { \R^2
} {}
adalah
\definitionsverweis {proyeksi stereografis}{}{.}

a) Tentukan
\mathl{\alpha_* \omega}{} untuk
\definitionsverweis {bentuk luas kanonik}{}{}
pada $S^2$.

b) Gunakan
\mathl{\alpha_* \omega}{} untuk menghitung luas permukaan bola.

}
{} {}

  \zwischenueberschrift{Soal untuk dikumpulkan}

\inputaufgabe
{5}
{

Kita meninjau
\definitionsverweis {grafik}{}{}
$M$ dari fungsi
\maabbeledisp {\psi} {\R^2} {\R
} {(x,y)} { y+x^2
} {,}
sebagai
\definitionsverweis {manifold Riemann}{}{.}
Hitung luas permukaan grafik di atas persegi
\mathl{[-1,1] \times [-1,1]}{.}

}
{} {}

\inputaufgabe
{5}
{

Misalkan

\mathdisp {M= \left\{ (x,x^2)  \mid   x \in \R        \right\}  \subseteq \R^2} {  }

adalah
\definitionsverweis {parabola}{}{,}
yaitu
\definitionsverweis {grafik}{}{}
dari
\definitionsverweis {fungsi}{}{}
\maabbeledisp {} {\R} {\R
} { x } { x^2
} {.}
Tunjukkan bahwa
\definitionsverweis {permukaan putar}{}{}
terkait yang diperoleh dengan memutarnya mengelilingi sumbu $x$ bukan suatu
\definitionsverweis {manifold}{}{.}

}
{} {}

\inputaufgabe
{4}
{

Nyatakan suatu
\definitionsverweis {permukaan bola}{}{} sebagai
\definitionsverweis {permukaan putar}{}{} dan gunakan representasi ini untuk menghitung luas permukaan bola tersebut.

}
{} {}

\inputaufgabe
{4}
{

Nyatakan suatu
\definitionsverweis {torus}{}{} sebagai
\definitionsverweis {permukaan putar}{}{} dan gunakan representasi ini untuk menghitung luas permukaannya.

}
{} {}

\inputaufgabe
{6}
{

Tentukan \anfuehrung{jarak}{} antara Osnabrück dan Bangalore
\zusatzklammer {gunakan $6370$ km sebagai jari-jari bumi} {} {}
dalam kedua pengertian berikut.

a) Menyusuri permukaan bumi
\zusatzklammer {jarak garis udara} {} {.}

b) Menembus bumi
\zusatzklammer {jarak garis tikus tanah} {} {.}

}
{} {}

\inputaufgabe
{6}
{

Berapakah panjang citra
\definitionsverweis {lingkaran lintang}{}{} pada lintang $30^\circ$
pada
\definitionsverweis {peta-peta}{}{}
yang diuraikan dalam
Contoh 17.6,
Contoh 17.7,
dan
Contoh 17.8
\zusatzklammer {gunakan $6370$ km sebagai jari-jari bumi} {} {?}

}
{} {}
